\begin{tablalafrox}{Diccionario de la cuenta 7.2. Fuente: (Bases Administrativas, art. 17, pp. 12–13, y Formulario T-14, p. 63); (Bases Técnicas del caso, secciones 17.2–17.6, pp. 32–34).}{tab:T14-dic-7-2}
  {L{1.2cm} Y Y Y C{0.9cm} L{2.3cm}}
  {\cab{Código} & \cab{Paquete de trabajo} & \cab{Entregable} & \cab{Criterio de aceptación} & \cab{Resp.} & \cab{Período}}
  7.2.1 & Innovación 2: informe de medición inicial y contrato de evidencia (INN-02.P1) & Informe de medición inicial y contrato de evidencia. & Aprobado por el Encargado de Seguridad. & CAL & E1 · Meses 2 y 3 \\
  7.2.2 & Innovación 2: circuito mínimo de reproducción en la preventa (INN-02.P2) & Circuito mínimo de reproducción en la preventa. & Un incidente se reproduce en QA a partir de su evidencia. & CAL & E1 · Meses 3 a 5 \\
  7.2.3 & Innovación 2: captura de evidencia en las aplicaciones de preventa, preparación y reparto (INN-02.P3) & Captura de evidencia en las aplicaciones de preventa, preparación y reparto. & Las tres aplicaciones envían evidencia válida antes del H4. & CAL & E1 · Meses 5 a 10 (antes del H4) \\
  7.2.4 & Innovación 2: informe de validación comparada durante la implementación (INN-02.P4, primera parte) & Informe de validación comparada durante la implementación. & Al mes 15 se cumplen dos metas: el tiempo de diagnóstico baja 30 \%, con al menos 20 casos comparables, y al menos el 80 \% de los incidentes se reproduce. & CAL & E1 · Meses 11 a 15 (validación en meses 13 a 15) \\
  7.2.5 & Innovación 2: mantenimiento de escenarios de regresión (INN-02.P4, segunda parte) & Informe trimestral. & Ningún incidente se reabre por la misma causa, y ninguna evidencia contiene datos personales. & CAL & Operación · Desde el mes 21. \\
\end{tablalafrox}
